% claim.tex -- AI4Math claim of record (AI-generated, 2026-09-30; instantiated
% from harness/templates/claim/claim.tex per SOP 08b).
% Machine-readable theorem anchors are the "% LEAN:" lines; scripts/paper-lint.sh
% and the claim audit check them. All Lean listings are verbatim, byte-identical
% contiguous excerpts of frozen artifacts of tasks/20260928-holestrip-tour-pipeline:
%   statements : formalized/holestrip-tour.statement.txt (84 lines, LF;
%                sha256 96875db4a784ff7a3862e3fe255c9e25c09f98a68d9d98d30700d90b4fd6929e;
%                byte-identical to formalized/holestrip-tour.lean)
%   proof file : attempts/stage4b-full.lean (130,382 bytes; 2,633 newline-terminated
%                lines, CRLF-dominant mixed endings on the working disk; excerpts
%                LF-normalised per the chorded-cycle precedent, byte-exact
%                original shipped in proofs/)
% Line ranges used: statement 43-76, 79-83; final 1589-1644, 2001-2036,
% 2447-2575, 2577-2633. The frozen statement file's literal "sorry"
% occurrences (the two theorem proof placeholders at lines 77 and 84, plus
% one mention in the head comment block at line 18) all lie outside the
% excerpt ranges, as required by the paper-lint listing gate; the final proof
% file contains zero sorry/admit/native_decide anywhere. Verification:
% audit/listing-verify.txt (audit/inspect-tex.py --verify).
\documentclass[11pt]{article}

\usepackage[T1]{fontenc}
\usepackage[utf8]{inputenc}
\usepackage{amsmath,amssymb,amsthm}
\usepackage{graphicx}
\usepackage{hyperref}
\usepackage{xurl}
\setlength{\emergencystretch}{3em}
% lstlean.tex — Lean style for the listings package.
% Lineage: Jeremy Avigad (2015), adapted from Assia Mahboubi's SSR style;
% inspired by Olivier Verdier's unixode.sty for the Unicode literate table.
% No CTAN package or upstream repo exists; papers carry this file in their
% source package (cap set arXiv:1907.01449, sphere eversion arXiv:2210.07746).
% AI4Math adaptation (AI-generated, 2026-09-26): sorry/admit/sorryAx elevated
% to a dedicated error tier, matching the pipeline's 0-sorry constitution.
\usepackage{listings}
\usepackage{xcolor}

\definecolor{keywordcolor}{rgb}{0.7, 0.1, 0.1}   % red keywords
\definecolor{tacticcolor}{rgb}{0.0, 0.1, 0.6}    % blue tactics
\definecolor{commentcolor}{rgb}{0.4, 0.4, 0.4}   % grey comments
\definecolor{symbolcolor}{rgb}{0.0, 0.1, 0.6}    % blue symbols
\definecolor{sortcolor}{rgb}{0.1, 0.5, 0.1}      % green sorts
\definecolor{errorcolor}{rgb}{0.6, 0, 0}         % dark red: sorry/admit
\definecolor{stringcolor}{rgb}{0.5, 0.1, 0.5}    % purple strings

\lstdefinelanguage{lean}{
  % Anything between $ becomes LaTeX math mode
  mathescape=false,
  % Comments may or not include Latex commands
  texcl=false,
  % keywords, list taken from lean-syntax.txt
  morekeywords=[1]{import, prelude, protected, private, noncomputable,
    definition, meta, renaming, hiding, parameter, parameters, begin,
    constant, constants, lemma, variable, variables, theory, print,
    theorem, example, open, section, namespace, universe, universes, end,
    modifier, modifiers, using, export, exposing, axiom, inductive,
    coinductive, with, structure, class, instance, attribute, local, where,
    by, have, show, let, in, if, then, else, match, do, for, fun, assume,
    from, take, calc, include, omit, infix, infixl, infixr, notation,
    macro, syntax, elab, set_option, initialize, deriving, opaque, abbrev,
    mutual, partial, scoped, termination_by, decreasing_by},
  % Sorts
  morekeywords=[2]{Sort, Type, Prop},
  % tactics, list taken from lean-syntax.txt (tactic keywords)
  morekeywords=[3]{intro, intros, apply, exact, refine, rfl, simp, simp_all,
    simpa, simp_rw, rw, rewrite, erewrite, ring, ring_nf, omega, decide,
    norm_num, norm_cast, induction, cases, rcases, obtain, constructor,
    ext, funext, conv, congr, field_simp, linarith, nlinarith, positivity,
    tauto, trivial, aesop, use, by_cases, by_contra, push_neg, contrapose,
    symm, trans, assumption, contradiction, exfalso, specialize,
    generalize, revert, clear, rename_i, subst, unfold, delta, change,
    convert, split_ifs, interval_cases, fin_cases, zify, gcongr,
    nth_rewrite, suffices, generalizing, at, only, native_decide},
  % warnings - constitution tier: must never survive into a final proof
  morekeywords=[4]{sorry, admit, sorryAx},
  literate=
    {α}{{\ensuremath{\alpha}}}1
    {β}{{\ensuremath{\beta}}}1
    {γ}{{\ensuremath{\gamma}}}1
    {δ}{{\ensuremath{\delta}}}1
    {ε}{{\ensuremath{\varepsilon}}}1
    {ζ}{{\ensuremath{\zeta}}}1
    {η}{{\ensuremath{\eta}}}1
    {θ}{{\ensuremath{\theta}}}1
    {ι}{{\ensuremath{\iota}}}1
    {κ}{{\ensuremath{\kappa}}}1
    {λ}{{\ensuremath{\lambda}}}1
    {μ}{{\ensuremath{\mu}}}1
    {ν}{{\ensuremath{\nu}}}1
    {ξ}{{\ensuremath{\xi}}}1
    {π}{{\ensuremath{\pi}}}1
    {ρ}{{\ensuremath{\rho}}}1
    {σ}{{\ensuremath{\sigma}}}1
    {τ}{{\ensuremath{\tau}}}1
    {υ}{{\ensuremath{\upsilon}}}1
    {φ}{{\ensuremath{\varphi}}}1
    {χ}{{\ensuremath{\chi}}}1
    {ψ}{{\ensuremath{\psi}}}1
    {ω}{{\ensuremath{\omega}}}1
    {Γ}{{\ensuremath{\Gamma}}}1
    {Δ}{{\ensuremath{\Delta}}}1
    {Θ}{{\ensuremath{\Theta}}}1
    {Λ}{{\ensuremath{\Lambda}}}1
    {Ξ}{{\ensuremath{\Xi}}}1
    {Π}{{\ensuremath{\Pi}}}1
    {Σ}{{\ensuremath{\Sigma}}}1
    {Φ}{{\ensuremath{\Phi}}}1
    {Ψ}{{\ensuremath{\Psi}}}1
    {Ω}{{\ensuremath{\Omega}}}1
    {ℵ}{{\ensuremath{\aleph}}}1
    {∀}{{\ensuremath{\forall}}}1
    {∃!}{{\ensuremath{\exists}!}}1
    {∃}{{\ensuremath{\exists}}}1
    {∈}{{\ensuremath{\in}}}1
    {∉}{{\ensuremath{\notin}}}1
    {∘}{{\ensuremath{\circ}}}1
    {∩}{{\ensuremath{\cap}}}1
    {∪}{{\ensuremath{\cup}}}1
    {⊆}{{\ensuremath{\subseteq}}}1
    {⊂}{{\ensuremath{\subset}}}1
    {⊇}{{\ensuremath{\supseteq}}}1
    {⊃}{{\ensuremath{\supset}}}1
    {⊄}{{\ensuremath{\nsubseteq}}}1
    {∅}{{\ensuremath{\emptyset}}}1
    {∧}{{\ensuremath{\wedge}}}1
    {∨}{{\ensuremath{\vee}}}1
    {¬}{{\ensuremath{\neg}}}1
    {⊤}{{\ensuremath{\top}}}1
    {⊥}{{\ensuremath{\bot}}}1
    {↔}{{\ensuremath{\leftrightarrow}}}1
    {→}{{\ensuremath{\rightarrow}}}1
    {←}{{\ensuremath{\leftarrow}}}1
    {↑}{{\ensuremath{\uparrow}}}1
    {⇒}{{\ensuremath{\Rightarrow}}}1
    {⇐}{{\ensuremath{\Leftarrow}}}1
    {⟹}{{\ensuremath{\Longrightarrow}}}1
    {↦}{{\ensuremath{\mapsto}}}1
    {≤}{{\ensuremath{\leq}}}1
    {≥}{{\ensuremath{\geq}}}1
    {≠}{{\ensuremath{\neq}}}1
    {≈}{{\ensuremath{\approx}}}1
    {≡}{{\ensuremath{\equiv}}}1
    {≃}{{\ensuremath{\simeq}}}1
    {≅}{{\ensuremath{\cong}}}1
    {×}{{\ensuremath{\times}}}1
    {·}{{\ensuremath{\cdot}}}1
    {±}{{\ensuremath{\pm}}}1
    {⊕}{{\ensuremath{\oplus}}}1
    {⊗}{{\ensuremath{\otimes}}}1
    {⋆}{{\ensuremath{\star}}}1
    {▸}{{\ensuremath{\blacktriangleright}}}1
    {⟨}{{\ensuremath{\langle}}}1
    {⟩}{{\ensuremath{\rangle}}}1
    {⦃}{{\ensuremath{\{\!\mid}}}1
    {⦄}{{\ensuremath{\mid\!\}}}}1
    {⟦}{{\ensuremath{[\![}}}1
    {⟧}{{\ensuremath{]\!]}}}1
    {⌊}{{\ensuremath{\lfloor}}}1
    {⌋}{{\ensuremath{\rfloor}}}1
    {⌈}{{\ensuremath{\lceil}}}1
    {⌉}{{\ensuremath{\rceil}}}1
    {√}{{\ensuremath{\surd}}}1
    {∣}{{\ensuremath{\mid}}}1
    {∤}{{\ensuremath{\nmid}}}1
    {∎}{{\ensuremath{\blacksquare}}}1
    {⋯}{{\ensuremath{\cdots}}}1
    {…}{{\ensuremath{\ldots}}}1
    {∑}{{\ensuremath{\sum}}}1
    {∏}{{\ensuremath{\prod}}}1
    {⁻¹}{{\ensuremath{^{-1}}}}1
    {¹}{{\ensuremath{^1}}}1
    {²}{{\ensuremath{^2}}}1
    {³}{{\ensuremath{^3}}}1
    {ⁿ}{{\ensuremath{^n}}}1
    {ᵐ}{{\ensuremath{^m}}}1
    {₀}{{\ensuremath{_0}}}1
    {₁}{{\ensuremath{_1}}}1
    {₂}{{\ensuremath{_2}}}1
    {₃}{{\ensuremath{_3}}}1
    {₄}{{\ensuremath{_4}}}1
    {₅}{{\ensuremath{_5}}}1
    {ₙ}{{\ensuremath{_n}}}1
    {ₘ}{{\ensuremath{_m}}}1
    {ₖ}{{\ensuremath{_k}}}1
    {ᵢ}{{\ensuremath{_i}}}1
    {ⱼ}{{\ensuremath{_j}}}1
    {ℤ}{{\ensuremath{\mathbb{Z}}}}1
    {ℕ}{{\ensuremath{\mathbb{N}}}}1
    {ℚ}{{\ensuremath{\mathbb{Q}}}}1
    {ℝ}{{\ensuremath{\mathbb{R}}}}1
    {ℂ}{{\ensuremath{\mathbb{C}}}}1
    {𝔽}{{\ensuremath{\mathbb{F}}}}1
    {𝒫}{{\ensuremath{\mathcal{P}}}}1
    {∗}{{\ensuremath{\ast}}}1
    {•}{{\ensuremath{\bullet}}}1,
  % Comments
  morecomment=[l]{--},
  morecomment=[s]{/-}{-/},
  morestring=[b]",
  stringstyle=\color{stringcolor},
  keepspaces=true,
  keywordstyle=[1]\color{keywordcolor}\bfseries,
  keywordstyle=[2]\color{sortcolor}\bfseries,
  keywordstyle=[3]\color{tacticcolor}\bfseries,
  keywordstyle=[4]\color{errorcolor}\bfseries\underline,
  escapeinside={(*@}{@*)},
  columns=fullflexible,
  basicstyle=\ttfamily\small,
  commentstyle=\color{commentcolor}\itshape,
  breaklines=true,
  showstringspaces=false,
  frame=single,
  framesep=2mm,
  xleftmargin=4mm,
  numbers=left,
  numberstyle=\tiny\color{commentcolor},
  numbersep=6pt,
}

% Inline Lean code in prose: \lean{Int.ModEq.pow}
\newcommand{\lean}[1]{\lstinline[language=lean,basicstyle=\ttfamily\normalsize]|#1|}

\lstset{language=lean}

% --- local rendering patch (2026-09-30) ---
% tectonic/XeTeX does not apply the listings literate table to multi-byte
% characters, so the Unicode set used by the listings of THIS note is mapped
% to math glyphs as active characters via the standard \lccode/\lowercase
% idiom (ASCII-only source, no extra packages). Characters deliberately NOT
% mapped: U+00B7 (middle dot) and U+2014 (em-dash), which xeCJK sets up at
% load time (making them active first aborts the run; observed 2026-09-27).
% U+2460-2463 (circled digits 1-4, case labels) ARE mapped, via the text-mode
% \textcircled macro (observed 2026-09-30: xeCJK does not carry them inside a
% monospaced listing). Rendering only: listing source bytes are
% untouched, so the byte-exactness checks are unaffected, and the \ifdefined
% guard leaves the pdflatex path identical to the template.
\ifdefined\XeTeXversion
  {\lccode`\~="00AC \lowercase{\global\catcode`~=\active \global\def~{\ensuremath{\neg}}}}
  {\lccode`\~="00D7 \lowercase{\global\catcode`~=\active \global\def~{\ensuremath{\times}}}}
  {\lccode`\~="0394 \lowercase{\global\catcode`~=\active \global\def~{\ensuremath{\Delta}}}}
  {\lccode`\~="03A3 \lowercase{\global\catcode`~=\active \global\def~{\ensuremath{\Sigma}}}}
  {\lccode`\~="03B1 \lowercase{\global\catcode`~=\active \global\def~{\ensuremath{\alpha}}}}
  {\lccode`\~="2115 \lowercase{\global\catcode`~=\active \global\def~{\ensuremath{\mathbb{N}}}}}
  {\lccode`\~="2190 \lowercase{\global\catcode`~=\active \global\def~{\ensuremath{\leftarrow}}}}
  {\lccode`\~="2192 \lowercase{\global\catcode`~=\active \global\def~{\ensuremath{\rightarrow}}}}
  {\lccode`\~="2194 \lowercase{\global\catcode`~=\active \global\def~{\ensuremath{\leftrightarrow}}}}
  {\lccode`\~="21A6 \lowercase{\global\catcode`~=\active \global\def~{\ensuremath{\mapsto}}}}
  {\lccode`\~="21D2 \lowercase{\global\catcode`~=\active \global\def~{\ensuremath{\Rightarrow}}}}
  {\lccode`\~="2200 \lowercase{\global\catcode`~=\active \global\def~{\ensuremath{\forall}}}}
  {\lccode`\~="2203 \lowercase{\global\catcode`~=\active \global\def~{\ensuremath{\exists}}}}
  {\lccode`\~="2208 \lowercase{\global\catcode`~=\active \global\def~{\ensuremath{\in}}}}
  {\lccode`\~="2209 \lowercase{\global\catcode`~=\active \global\def~{\ensuremath{\notin}}}}
  {\lccode`\~="2212 \lowercase{\global\catcode`~=\active \global\def~{\ensuremath{-}}}}
  {\lccode`\~="2227 \lowercase{\global\catcode`~=\active \global\def~{\ensuremath{\wedge}}}}
  {\lccode`\~="2228 \lowercase{\global\catcode`~=\active \global\def~{\ensuremath{\vee}}}}
  {\lccode`\~="222A \lowercase{\global\catcode`~=\active \global\def~{\ensuremath{\cup}}}}
  {\lccode`\~="2260 \lowercase{\global\catcode`~=\active \global\def~{\ensuremath{\neq}}}}
  {\lccode`\~="2264 \lowercase{\global\catcode`~=\active \global\def~{\ensuremath{\leq}}}}
  {\lccode`\~="2265 \lowercase{\global\catcode`~=\active \global\def~{\ensuremath{\geq}}}}
  {\lccode`\~="2286 \lowercase{\global\catcode`~=\active \global\def~{\ensuremath{\subseteq}}}}
  {\lccode`\~="2460 \lowercase{\global\catcode`~=\active \global\def~{\textcircled{\tiny 1}}}}
  {\lccode`\~="2461 \lowercase{\global\catcode`~=\active \global\def~{\textcircled{\tiny 2}}}}
  {\lccode`\~="2462 \lowercase{\global\catcode`~=\active \global\def~{\textcircled{\tiny 3}}}}
  {\lccode`\~="2463 \lowercase{\global\catcode`~=\active \global\def~{\textcircled{\tiny 4}}}}
  {\lccode`\~="27E8 \lowercase{\global\catcode`~=\active \global\def~{\ensuremath{\langle}}}}
  {\lccode`\~="27E9 \lowercase{\global\catcode`~=\active \global\def~{\ensuremath{\rangle}}}}
  {\lccode`\~="27FA \lowercase{\global\catcode`~=\active \global\def~{\ensuremath{\Longleftrightarrow}}}}
\fi
\ifdefined\XeTeXversion
  \usepackage{xeCJK}
  \setCJKmainfont{SimSun}
\fi

\newtheorem{theorem}{Theorem}

\title{The 3\,$\times$\,$n$ holed-strip knight graph with two far corners
removed: open and closed tour existence characterisation\\ (Lean 4
machine-checked claim of record)}
\author{Aurora0134\thanks{Contact: via the GitHub account Aurora0134; ORCID
placeholder --- to be filled in by the author before any upload. No personal
information is carried in this note.}}
\date{September 30, 2026}

\begin{document}
\maketitle

\begin{abstract}
Let $K_n$ be the knight graph of the $3\times n$ chessboard with the two far
corners $(0,0)$ and $(2,n-1)$ removed. We prove, with machine-checked Lean 4
proofs, the exact existence characterisation of tours on $K_n$: an open tour
(Hamiltonian path) exists if and only if $n\in\{1,4,5\}\cup\{n\ge 8\}$, and a
closed tour (Hamiltonian cycle) exists if and only if $n$ is even and
$n\ge 8$. Both theorems compile with no \texttt{sorry}, and
\texttt{\#print axioms} reports only \{propext, Quot.sound,
Classical.choice\}. Four rounds of novelty review found no occupying record
for the family: the nearest published neighbours remove one square from
$3\times n$ boards or two squares from $4\times n$ boards, and none of them
treats two-square $3\times n$ boards. This note is a priority claim of
record, not a full paper.
\end{abstract}

\section{Introduction}

We study the knight graph $K_n$ of the $3\times n$ chessboard with the two
far corners removed --- a deficient-board variant of the classical $3\times
n$ knight's-tour problem whose undirected tour counts are registered as OEIS
A169696 \cite{a169696}. For this variant we establish the complete
existence characterisation for both open and closed tours, formalised and
proved in Lean 4 with every assertion adjudicated by the kernel. The result
was produced by the AI4Math pipeline (LLM generation under Lean-kernel
adjudication); this note is a priority claim of record, and the full
narrative paper is prepared separately. All formal statements and proofs
were checked by the Lean 4 kernel: every proof compiles with no
\texttt{sorry}, and \texttt{\#print axioms} reports only \{propext,
Quot.sound, Classical.choice\}.

\section{Statement}

Throughout, a board cell is a pair $(r,c)$ with $r\in\{0,1,2\}$ and
$c\in\{0,\dots,n-1\}$; $K_n$ has as vertices the $3n-2$ cells other than
$(0,0)$ and $(2,n-1)$, two vertices being adjacent when they are a knight's
move apart (row difference $1$ and column difference $2$, or row difference
$2$ and column difference $1$). An open tour is a Hamiltonian path of
$K_n$; a closed tour is a Hamiltonian cycle. Both notions use the standard
mathlib predicates \cite{mathlib} (\texttt{Walk.IsHamiltonian} and
\texttt{Walk.IsHamiltonianCycle}).

% LEAN: tasks/20260928-holestrip-tour-pipeline :: exists_hamiltonian_path_iff  grade=完全证明
\begin{theorem}[open tours]\label{thm:open}
$K_n$ has an open tour if and only if $n=1$, or $n=4$, or $n=5$, or
$n\ge 8$.
\end{theorem}

% LEAN: tasks/20260928-holestrip-tour-pipeline :: exists_hamiltonian_cycle_iff  grade=完全证明
\begin{theorem}[closed tours]\label{thm:closed}
$K_n$ has a closed tour if and only if $n$ is even and $n\ge 8$.
\end{theorem}

The case $n=1$ in Theorem~\ref{thm:open} is the one-point graph: the empty
walk is a (trivial) Hamiltonian path. This reading was fixed before the
statement freeze (the pre-freeze correction F0b recorded in the task card,
backed by a kernel probe); the closed-tour side is unaffected because
\texttt{IsHamiltonianCycle} excludes the empty walk.

\begin{lstlisting}[caption={Frozen statement (verbatim lines 43--76 of the
frozen statement file, sha256 \texttt{96875db4...}): board encoding, the
knight graph, and the head of Theorem~\ref{thm:open}. The proof body
placeholder lies outside the excerpt.},label={lst:stmt-open}]
/-- 3×n 棋盘的格子坐标。 -/
abbrev Cell (n : ℕ) := Fin 3 × Fin n

/-- 棋盘格子（非洞）：删去 (0,0) 与 (2, n−1)，用 `.val` 编码（见文件头注释）。 -/
abbrev IsCell (n : ℕ) (p : Cell n) : Prop :=
  ¬ (p.1.val = 0 ∧ p.2.val = 0) ∧ ¬ (p.1.val = 2 ∧ p.2.val + 1 = n)

/-- 顶点类型：3n−2 个非洞格子。 -/
def V (n : ℕ) := {p : Cell n // IsCell n p}

instance (n : ℕ) : Fintype (V n) := Subtype.fintype _
instance (n : ℕ) : DecidableEq (V n) := Subtype.instDecidableEq

/-- 骑士走法的四个代表方向（b − a ∈ {(1,2),(2,1),(−1,2),(−2,1)}）；
    `fromRel` 自动对称化，补全另外四个方向（(|Δ行|,|Δ列|) ∈ {(1,2),(2,1)} 全 8 组合）。 -/
abbrev leap {n : ℕ} (a b : Cell n) : Prop :=
  ((a.1 : ℕ) + 1 = (b.1 : ℕ) ∧ (a.2 : ℕ) + 2 = (b.2 : ℕ)) ∨
  ((a.1 : ℕ) + 2 = (b.1 : ℕ) ∧ (a.2 : ℕ) + 1 = (b.2 : ℕ)) ∨
  ((b.1 : ℕ) + 1 = (a.1 : ℕ) ∧ (a.2 : ℕ) + 2 = (b.2 : ℕ)) ∨
  ((b.1 : ℕ) + 2 = (a.1 : ℕ) ∧ (a.2 : ℕ) + 1 = (b.2 : ℕ))

/-- 完整 8 向（对称）骑士关系。 -/
abbrev leap8 {n : ℕ} (a b : Cell n) : Prop := leap a b ∨ leap b a

/-- K_n：3×n 棋盘删 (0,0)、(2,n−1) 后的骑士图。 -/
def knightGraph (n : ℕ) : SimpleGraph (V n) :=
  SimpleGraph.fromRel (fun a b : V n => leap8 a.1 b.1)

/-- **定理 O**（开巡游，命题 (a) 前半）：K_n 存在哈密顿路径
    当且仅当 n ∈ {4, 5} ∪ {n ≥ 8}。
    n=1 为单点图平凡路径情形（见文件头修正注记）。 -/
theorem exists_hamiltonian_path_iff (n : ℕ) :
    (∃ a b : V n, ∃ w : (knightGraph n).Walk a b, w.IsHamiltonian) ↔
      n = 1 ∨ n = 4 ∨ n = 5 ∨ 8 ≤ n := by
\end{lstlisting}

\begin{lstlisting}[caption={Frozen statement, verbatim lines 79--83: the head
of Theorem~\ref{thm:closed}.},label={lst:stmt-closed}]
/-- **定理 C**（闭巡游，命题 (a) 后半）：K_n 存在哈密顿圈
    当且仅当 n 为偶数且 n ≥ 8。 -/
theorem exists_hamiltonian_cycle_iff (n : ℕ) :
    (∃ a : V n, ∃ w : (knightGraph n).Walk a a, w.IsHamiltonianCycle) ↔
      Even n ∧ 8 ≤ n := by
\end{lstlisting}

\section{Proof}

The complete machine proof is the file \texttt{stage4b-full.lean} (2{,}633
lines), shipped byte-exactly in \texttt{proofs/}; the listings below are
verbatim contiguous excerpts, in the order in which they appear in that
file. In outline: the positive side is an induction machine on explicit
vertex lists --- closed tours for even $n\ge 8$ by four explicit base lists
($n=8,10,12,14$) plus an insertion step that splices a fixed 8-column wedge
into a designated break edge (strong induction over $n \bmod 8$); open tours
for odd $n\ge 9$ by two base lists ($n=9,11$) plus a $4$-column step; the
small cases $n=1,4,5$ by explicit witness lists whose adjacency, coverage
and non-repetition are checked by \texttt{decide}. The negative side is
degree and parity arguments for $n\le 4$ and the odd closed case, plus ---
for the two hard zero values $n=6,7$ --- a cut barrier: a general lemma
stating that if deleting a vertex list $S$ leaves more than $|S|+1$ labelled
blocks, then no Hamiltonian path exists, applied to explicit certificates
whose correctness is decided by the kernel. The two final theorems assemble
these pieces by case analysis on $n$ (\texttt{omega} splits $n<8$ into
$0,\dots,7$).

\begin{lstlisting}[caption={Verbatim lines 1589--1644: the closed-tour
induction assembly --- \texttt{closed\_good\_list} (strong induction,
bases $8/10/12/14$, step $n\mapsto n+8$) and \texttt{closed\_cycle\_exists}
(loading the good list as a Hamiltonian cycle).},label={lst:proof-closed}]
/-- 闭侧好列表存在性：偶数 n ≥ 8 ⇒ 存在 GoodCycleList
    （强归纳：小 n 四基座 8/10/12/14，步进 m ↦ m+8）。 -/
theorem closed_good_list : ∀ n : ℕ, 8 ≤ n → Even n → ∃ l : List (V n), GoodCycleList l := by
  intro n
  apply Nat.strong_induction_on
    (p := fun n => 8 ≤ n → Even n → ∃ l : List (V n), GoodCycleList l)
  intro n ih hn heven
  rcases Nat.lt_or_ge n 16 with h16 | h16
  · -- 小 n：8 ≤ n < 16 且偶数 ⇒ n ∈ {8, 10, 12, 14}
    rcases heven with ⟨m, hm⟩
    have hn' : n = 8 ∨ n = 10 ∨ n = 12 ∨ n = 14 := by omega
    rcases hn' with h | h | h | h
    · subst h; exact ⟨baseV8, baseV8_good⟩
    · subst h; exact ⟨baseV10, baseV10_good⟩
    · subst h; exact ⟨baseV12, baseV12_good⟩
    · subst h; exact ⟨baseV14, baseV14_good⟩
  · -- n ≥ 16：从 n − 8 步进
    have hn8 : n - 8 < n := by omega
    have hn8' : 8 ≤ n - 8 := by omega
    have heven8 : Even (n - 8) := by
      rcases heven with ⟨m, hm⟩
      exact ⟨m - 4, by omega⟩
    obtain ⟨l, hg⟩ := ih (n - 8) hn8 hn8' heven8
    obtain ⟨l', hg'⟩ := step_closed_list hn8' hg
    rw [← show n - 8 + 8 = n from by omega]
    exact ⟨l', hg'⟩

/-- 闭巡游存在性（正向）：偶数 n ≥ 8 ⇒ K_n 存在哈密顿圈
    （GoodCycleList 经 cycleOfChain 装载为 IsHamiltonianCycle）。 -/
theorem closed_cycle_exists : ∀ n : ℕ, Even n → 8 ≤ n →
    ∃ a : V n, ∃ w : (knightGraph n).Walk a a, w.IsHamiltonianCycle := by
  intro n heven hn
  obtain ⟨l, hg⟩ := closed_good_list n hn heven
  obtain ⟨hne, hnd, hcov, hchain, hclose, j, hj1r, hj1c, hj2r, hj2c⟩ := hg
  obtain ⟨v1, v2, v3, h12, h13, h23⟩ := three_vertices_exist hn
  have hlen : 3 ≤ l.length := by
    have hsub : [v1, v2, v3] ⊆ l := fun a ha => by
      rw [List.mem_cons, List.mem_cons, List.mem_singleton] at ha
      rcases ha with h | h | h
      · rw [h]; exact hcov v1
      · rw [h]; exact hcov v2
      · rw [h]; exact hcov v3
    have hnd3 : ([v1, v2, v3] : List (V n)).Nodup := by
      rw [List.nodup_cons, List.nodup_cons]
      refine ⟨?_, ?_, List.nodup_singleton v3⟩
      · intro h
        rw [List.mem_cons, List.mem_singleton] at h
        rcases h with h | h
        · exact h12 h
        · exact h13 h
      · intro h
        rw [List.mem_singleton] at h
        exact h23 h
    have := (hnd3.subperm hsub).length_le
    simpa using this
  exact cycleOfChain l hne hnd hlen hcov hchain (hclose hne)
\end{lstlisting}

\begin{lstlisting}[caption={Verbatim lines 2001--2036: the open-tour
induction assembly for odd $n\ge 9$ (bases $9/11$, step $n\mapsto n+4$).},label={lst:proof-open}]
/-! ### 12.5 奇归纳组装（T3a 收官） -/

/-- 开侧好列表存在性：奇数 n ≥ 9 ⇒ 存在 GoodOpenList
    （强归纳：小 n 两基座 9/11，步进 m ↦ m+4）。 -/
theorem open_good_list : ∀ n : ℕ, 9 ≤ n → Odd n → ∃ l : List (V n), GoodOpenList l := by
  intro n
  apply Nat.strong_induction_on
    (p := fun n => 9 ≤ n → Odd n → ∃ l : List (V n), GoodOpenList l)
  intro n ih hn hodd
  rcases Nat.lt_or_ge n 13 with h13 | h13
  · -- 小 n：奇数且 9 ≤ n < 13 ⇒ n ∈ {9, 11}
    have hn' : n = 9 ∨ n = 11 := by
      rcases hodd with ⟨m, hm⟩
      omega
    rcases hn' with h | h
    · subst h; exact ⟨openBase9V, openBase9V_good⟩
    · subst h; exact ⟨openBase11V, openBase11V_good⟩
  · -- n ≥ 13：从 n − 4 步进
    have hn4 : n - 4 < n := by omega
    have hn4' : 9 ≤ n - 4 := by omega
    have hodd4 : Odd (n - 4) := by
      rcases hodd with ⟨m, hm⟩
      exact ⟨m - 2, by omega⟩
    obtain ⟨l, hg⟩ := ih (n - 4) hn4 hn4' hodd4
    obtain ⟨l', hg'⟩ := step_open_list hn4' hg
    rw [← show n - 4 + 4 = n from by omega]
    exact ⟨l', hg'⟩

/-- 开巡游存在性（正向·奇）：奇数 n ≥ 9 ⇒ K_n 存在哈密顿路径
    （GoodOpenList 经 pathOfChain 装载为 IsHamiltonian）。 -/
theorem open_tour_exists_odd : ∀ n : ℕ, 9 ≤ n → Odd n →
    ∃ a b : V n, ∃ w : (knightGraph n).Walk a b, w.IsHamiltonian := by
  intro n hn hodd
  obtain ⟨l, hg⟩ := open_good_list n hn hodd
  obtain ⟨hne, hnd, hcov, hchain, hend⟩ := hg
  exact pathOfChain l hne hnd hcov hchain
\end{lstlisting}

\begin{lstlisting}[caption={Verbatim lines 2447--2575: the cut barrier. The
generic lemmas \texttt{cut\_card\_bound} (label count along a chain),
\texttt{cut\_bound} and \texttt{no\_hamiltonian\_of\_cut}, then the explicit
$n=6$ certificate (6 middle-column vertices, 8 blocks) and $n=7$ certificate
(5 vertices, 7 blocks). The two \texttt{set\_option} resource-relaxation
lines preceding this section are omitted from the excerpt; the full file
carries them.},label={lst:proof-cut}]
/-! ## 15. 割集阻障（T4-open：n = 6/7 开侧否定） -/

theorem cut_card_bound {α : Type*} [DecidableEq α] {S : List α} {N : ℕ} {lab : α → Fin N}
    {G : SimpleGraph α}
    (hadj : ∀ a b : α, G.Adj a b → a ∉ S → b ∉ S → lab a = lab b) :
    ∀ {L : List α}, L.IsChain G.Adj →
      ((L.filter (· ∉ S)).map lab).toFinset.card ≤ 1 + L.tail.countP (· ∈ S) := by
  intro L
  induction L with
  | nil => intro _; simp
  | cons x xs ih =>
      intro hL
      cases xs with
      | nil => by_cases hx : x ∈ S <;> simp [hx]
      | cons y ys =>
          obtain ⟨hxy, htail⟩ := List.isChain_cons_cons.mp hL
          have hih := ih htail
          rw [List.tail_cons] at hih
          rw [List.tail_cons]
          by_cases hx : x ∈ S <;> by_cases hy : y ∈ S
          · -- ① x ∈ S, y ∈ S
            have hF : (x :: y :: ys).filter (· ∉ S) = (y :: ys).filter (· ∉ S) :=
              List.filter_cons_of_neg (by simp [hx])
            rw [hF]
            have hc : (y :: ys).countP (· ∈ S) = ys.countP (· ∈ S) + 1 := by simp [hy]
            rw [hc]; omega
          · -- ② x ∈ S, y ∉ S
            have hF : (x :: y :: ys).filter (· ∉ S) = (y :: ys).filter (· ∉ S) :=
              List.filter_cons_of_neg (by simp [hx])
            rw [hF]
            have hc : (y :: ys).countP (· ∈ S) = ys.countP (· ∈ S) := by simp [hy]
            rw [hc]; omega
          · -- ③ x ∉ S, y ∈ S
            have hF : (x :: y :: ys).filter (· ∉ S) = x :: (y :: ys).filter (· ∉ S) :=
              List.filter_cons_of_pos (by simp [hx])
            rw [hF, List.map_cons, List.toFinset_cons]
            have hc : (y :: ys).countP (· ∈ S) = ys.countP (· ∈ S) + 1 := by simp [hy]
            rw [hc]
            exact le_trans (Finset.card_insert_le (lab x) _) (by omega)
          · -- ④ x ∉ S, y ∉ S
            have hF : (x :: y :: ys).filter (· ∉ S) = x :: (y :: ys).filter (· ∉ S) :=
              List.filter_cons_of_pos (by simp [hx])
            rw [hF, List.map_cons, List.toFinset_cons]
            have hc : (y :: ys).countP (· ∈ S) = ys.countP (· ∈ S) := by simp [hy]
            rw [hc]
            have hlab : lab x = lab y := hadj x y hxy hx hy
            have hymem : lab y ∈ (((y :: ys).filter (· ∉ S)).map lab).toFinset :=
              List.mem_toFinset.mpr (List.mem_map_of_mem (f := lab)
                (List.mem_filter.mpr ⟨by simp, by simp [hy]⟩))
            rw [hlab, Finset.insert_eq_of_mem hymem]
            omega

theorem cut_bound {V : Type*} [Fintype V] [DecidableEq V] {G : SimpleGraph V} {N : ℕ}
    (S : List V) (hSnd : S.Nodup) (lab : V → Fin N)
    (hadj : ∀ a b : V, G.Adj a b → a ∉ S → b ∉ S → lab a = lab b)
    (hsurj : ∀ i : Fin N, ∃ v : V, v ∉ S ∧ lab v = i)
    {a b : V} {w : G.Walk a b} (hw : w.IsHamiltonian) :
    N ≤ S.length + 1 := by
  have hchain : w.support.IsChain G.Adj := SimpleGraph.Walk.isChain_adj_support w
  have hnd : w.support.Nodup :=
    List.nodup_iff_count_le_one.mpr (fun v => by have h := hw v; omega)
  have hcut := cut_card_bound hadj hchain
  have hcover : ∀ i : Fin N, i ∈ ((w.support.filter (· ∉ S)).map lab).toFinset := by
    intro i
    obtain ⟨v, hvS, hvl⟩ := hsurj i
    have hvm : v ∈ w.support := List.count_pos_iff.mp (by have h := hw v; omega)
    rw [← hvl]
    exact List.mem_toFinset.mpr (List.mem_map_of_mem (f := lab)
      (List.mem_filter.mpr ⟨hvm, by simp [hvS]⟩))
  have hNcard : N = (Finset.univ : Finset (Fin N)).card := by
    rw [Finset.card_univ, Fintype.card_fin]
  have hcle : (Finset.univ : Finset (Fin N)).card
      ≤ ((w.support.filter (· ∉ S)).map lab).toFinset.card :=
    Finset.card_le_card (fun i _ => hcover i)
  have htail : w.support.tail.countP (· ∈ S) ≤ w.support.countP (· ∈ S) :=
    List.Sublist.countP_le (List.tail_sublist w.support)
  have hlen : w.support.countP (· ∈ S) = (w.support.filter (· ∈ S)).length :=
    List.countP_eq_length_filter
  have hsub : List.Subperm (w.support.filter (· ∈ S)) S :=
    List.Nodup.subperm (hnd.filter _)
      (fun x hx => of_decide_eq_true (List.mem_filter.mp hx).2)
  have hflen := hsub.length_le
  omega

theorem no_hamiltonian_of_cut {V : Type*} [Fintype V] [DecidableEq V] {G : SimpleGraph V} {N : ℕ}
    (S : List V) (hSnd : S.Nodup) (lab : V → Fin N)
    (hadj : ∀ a b : V, G.Adj a b → a ∉ S → b ∉ S → lab a = lab b)
    (hsurj : ∀ i : Fin N, ∃ v : V, v ∉ S ∧ lab v = i)
    (hN : S.length + 1 < N) :
    ¬ ∃ a b : V, ∃ w : G.Walk a b, w.IsHamiltonian := by
  rintro ⟨a, b, w, hw⟩
  have h := cut_bound S hSnd lab hadj hsurj hw
  omega

/-- n=6 顶点构造助手。 -/
def v6 (r : Fin 3) (c : Fin 6) (h : IsCell 6 (r, c) := by decide) : V 6 := ⟨(r, c), h⟩

/-- n=6 割集证书：中间两列 6 格。 -/
def cutS6 : List (V 6) := [v6 0 2, v6 0 3, v6 1 2, v6 1 3, v6 2 2, v6 2 3]

/-- n=6 连通块标签（S 格与未列格 ↦ 0）。 -/
def cutLab6 : V 6 → Fin 8 := fun v =>
  match v.1.1.val, v.1.2.val with
  | 0, 1 => 0 | 0, 4 => 1 | 0, 5 => 2
  | 1, 0 => 3 | 1, 1 => 4 | 1, 4 => 5 | 1, 5 => 6
  | 2, 0 => 0 | 2, 1 => 7 | 2, 4 => 2
  | _, _ => 0

theorem no_open_tour_6 :
    ¬ ∃ a b : V 6, ∃ w : (knightGraph 6).Walk a b, w.IsHamiltonian :=
  no_hamiltonian_of_cut cutS6 (by decide) cutLab6 (by decide) (by decide) (by decide)

/-- n=7 顶点构造助手。 -/
def v7 (r : Fin 3) (c : Fin 7) (h : IsCell 7 (r, c) := by decide) : V 7 := ⟨(r, c), h⟩

/-- n=7 割集证书：5 格。 -/
def cutS7 : List (V 7) := [v7 0 2, v7 0 4, v7 1 3, v7 2 2, v7 2 4]

/-- n=7 连通块标签（S 格与未列格 ↦ 0）。 -/
def cutLab7 : V 7 → Fin 7 := fun v =>
  match v.1.1.val, v.1.2.val with
  | 0, 1 => 0 | 0, 3 => 1 | 0, 5 => 2 | 0, 6 => 3
  | 1, 0 => 4 | 1, 1 => 1 | 1, 2 => 0 | 1, 4 => 3 | 1, 5 => 1 | 1, 6 => 5
  | 2, 0 => 0 | 2, 1 => 6 | 2, 3 => 1 | 2, 5 => 3
  | _, _ => 0

theorem no_open_tour_7 :
    ¬ ∃ a b : V 7, ∃ w : (knightGraph 7).Walk a b, w.IsHamiltonian :=
  no_hamiltonian_of_cut cutS7 (by decide) cutLab7 (by decide) (by decide) (by decide)
\end{lstlisting}

\begin{lstlisting}[caption={Verbatim lines 2577--2633: the assembly of both
iff theorems from the stage lemmas, followed by the two closing
\texttt{\#print axioms} commands.},label={lst:proof-assembly}]
/-! ## 16. T5 组装：双 iff 主定理（statement 冻结件） -/

/-- **定理 O**（开巡游 iff）：K_n 存在哈密顿路径 ⟺ n = 1 ∨ 4 ∨ 5 ∨ ≥8。 -/
theorem exists_hamiltonian_path_iff (n : ℕ) :
    (∃ a b : V n, ∃ w : (knightGraph n).Walk a b, w.IsHamiltonian) ↔
      n = 1 ∨ n = 4 ∨ n = 5 ∨ 8 ≤ n := by
  constructor
  · intro h
    by_cases h8 : 8 ≤ n
    · exact Or.inr (Or.inr (Or.inr h8))
    · push_neg at h8
      have : n = 0 ∨ n = 1 ∨ n = 2 ∨ n = 3 ∨ n = 4 ∨ n = 5 ∨ n = 6 ∨ n = 7 := by omega
      rcases this with rfl | rfl | rfl | rfl | rfl | rfl | rfl | rfl
      · exact absurd h no_open_tour_0
      · exact Or.inl rfl
      · exact absurd h no_open_tour_2
      · exact absurd h no_open_tour_3
      · exact Or.inr (Or.inl rfl)
      · exact Or.inr (Or.inr (Or.inl rfl))
      · exact absurd h no_open_tour_6
      · exact absurd h no_open_tour_7
  · rintro (rfl | rfl | rfl | h8)
    · exact open_tour_1
    · exact open_tour_4
    · exact open_tour_5
    · rcases Nat.even_or_odd n with he | ho
      · exact open_of_closed (closed_cycle_exists n he h8)
      · obtain ⟨k, rfl⟩ := ho
        exact open_tour_exists_odd _ (by omega) ⟨k, rfl⟩

/-- **定理 C**（闭巡游 iff）：K_n 存在哈密顿圈 ⟺ Even n ∧ ≥8。 -/
theorem exists_hamiltonian_cycle_iff (n : ℕ) :
    (∃ a : V n, ∃ w : (knightGraph n).Walk a a, w.IsHamiltonianCycle) ↔
      Even n ∧ 8 ≤ n := by
  constructor
  · intro h
    constructor
    · rcases Nat.even_or_odd n with he | ho
      · exact he
      · exact absurd h (no_cycle_odd ho)
    · by_contra h8
      push_neg at h8
      have : n = 0 ∨ n = 1 ∨ n = 2 ∨ n = 3 ∨ n = 4 ∨ n = 5 ∨ n = 6 ∨ n = 7 := by omega
      rcases this with rfl | rfl | rfl | rfl | rfl | rfl | rfl | rfl
      · exact no_closed_tour_0 h
      · exact no_cycle_odd ⟨0, rfl⟩ h
      · exact no_closed_tour_2 h
      · exact no_cycle_odd ⟨1, rfl⟩ h
      · exact no_closed_tour_4 h
      · exact no_cycle_odd ⟨2, rfl⟩ h
      · exact no_closed_tour_6 h
      · exact no_cycle_odd ⟨3, rfl⟩ h
  · rintro ⟨he, h8⟩
    exact closed_cycle_exists n he h8

#print axioms exists_hamiltonian_path_iff
#print axioms exists_hamiltonian_cycle_iff
\end{lstlisting}

\section{Verification evidence}

\begin{itemize}
\item Toolchain: Lean \texttt{v4.34.0} (\texttt{leanprover/lean4:v4.34.0}),
mathlib4 \texttt{v4.34.0}, revision \texttt{5ed29652}. Reproduction: with
that pin, \texttt{lake env lean} on
\texttt{proofs/01-holestrip-tour-proved.lean} compiles with exit code $0$.
\item Statement integrity: the frozen statement file (84 lines) has sha256
\begin{center}
\texttt{96875db4a784ff7a3862e3fe255c9e25c09f98}\allowbreak
\texttt{a68d9d98d30700d90b4fd6929e},
\end{center}
byte-identical to the frozen Lean file \texttt{holestrip-tour.lean}; the
final proof file (130{,}382 bytes) has sha256
\begin{center}
\texttt{9cf0b5b0af8cedb2edfd9d63779a2e8d4a6c973}\allowbreak
\texttt{237b616c1482a05bbf1e0a70b}.
\end{center}
A re-run of the audit department's symmetric-difference method (comments
stripped, declarations split, definitions compared full-block, theorems up
to \texttt{:=}) reports all eight frozen declarations header-identical
(six definitions and the two main theorems), with no frozen-only
declaration and proof-body scaffolding added on the final side
(\texttt{audit/statement-diff.out}).
\item Kernel check: compiles with $0$ \texttt{sorry} (whole-file scan
including comments, also covering \texttt{admit} and \texttt{native\_decide}:
zero hits); \texttt{\#print axioms} output, verbatim (kernel re-run of
2026-09-30, exit code 0, kept in \texttt{audit/axioms.out}):
\begin{lstlisting}[basicstyle=\ttfamily\footnotesize,frame=none]
'open_tour_exists_odd' depends on axioms: [propext, Classical.choice, Quot.sound]
'open_good_list' depends on axioms: [propext, Classical.choice, Quot.sound]
'step_open_list' depends on axioms: [propext, Classical.choice, Quot.sound]
'closed_cycle_exists' depends on axioms: [propext, Classical.choice, Quot.sound]
'openBase9V_good' depends on axioms: [propext, Classical.choice, Quot.sound]
'openBase11V_good' depends on axioms: [propext, Classical.choice, Quot.sound]
'card_V' depends on axioms: [propext, Classical.choice, Quot.sound]
'no_cycle_odd' depends on axioms: [propext, Classical.choice, Quot.sound]
'no_open_tour_2' depends on axioms: [propext, Classical.choice, Quot.sound]
'no_closed_tour_2' depends on axioms: [propext, Classical.choice, Quot.sound]
'no_open_tour_3' depends on axioms: [propext, Classical.choice, Quot.sound]
'no_closed_tour_3' depends on axioms: [propext, Classical.choice, Quot.sound]
'open_tour_1' depends on axioms: [propext, Classical.choice, Quot.sound]
'no_open_tour_0' depends on axioms: [propext, Quot.sound]
'no_closed_tour_0' depends on axioms: [propext, Quot.sound]
'no_closed_tour_4' depends on axioms: [propext, Classical.choice, Quot.sound]
'no_closed_tour_6' depends on axioms: [propext, Classical.choice, Quot.sound]
'exists_hamiltonian_path_iff' depends on axioms: [propext, Classical.choice, Quot.sound]
'exists_hamiltonian_cycle_iff' depends on axioms: [propext, Classical.choice, Quot.sound]
\end{lstlisting}
The two main theorems depend on exactly \{propext, Classical.choice,
Quot.sound\}; every line is a subset of the whitelist.
\item Grading by the audit department: both proven theorems
\emph{complete proof}; the lowest grade reached for the case = \emph{complete
proof} (final adjudication \texttt{gate3-final-audit.md}: statement
symmetric difference $8/8$ identical, strict compile exit $0$, zero
\texttt{sorry}/\texttt{admit}, axiom whitelist subset, no added hypotheses).
The adjudication also records the non-blocking disclosures carried into
Section~\ref{sec:limits}.
\item Audit chain (originals in this deposit under \texttt{audit/} unless
noted): the final adjudication, the symmetric-difference re-run, the axiom
re-print, and this note's claim check \texttt{claim-check.md} --- added to
\texttt{audit/} after the audit department returns its adjudication, and a
prerequisite for publication.
\end{itemize}

\section{Novelty evidence}\label{sec:novelty}

Adjudication copied from the pipeline's exclusivity reviews (C9 gate), four
rounds of which ran before the problem was accepted and during it: verdict
\textbf{no occupying record found} for the existence characterisation of
open and closed tours on $K_n$. Value tier: a small original classification
theorem. The nearest published neighbours are delimited explicitly and none
of them is this family: Miller--Farnsworth \cite{millerfarnsworth} classify
closed tours on $3\times n$ boards with \emph{one} square removed;
Srichote et al.\ \cite{srichote} classify closed tours on $4\times n$ boards
with \emph{two} squares removed; Bi et al.\ \cite{bi} raised the two-square
question for $4\times n$ only; DeMaio--Hippchen \cite{demaio} give the
minimal number of squares to remove, not positions; Schwenk
\cite{schwenk} and the classical line treat full boards. The registered
counting sequence for full $3\times n$ boards, OEIS A169696 \cite{a169696},
is a different object (no holes; counting, not existence). The full query
log, including every zero-hit query, is in this deposit at
\texttt{audit/c9-record.md} (four source records, byte-for-byte).

\begin{center}
\resizebox{\textwidth}{!}{%
\begin{tabular}{lll}
\hline
channel & scope & result \\
\hline
literature (journals/theses/preprints) & nearest neighbours read at
full-text level (Miller--Farnsworth 2013; Srichote et al.\ 2022; Bi et
al.\ 2015; DeMaio--Hippchen 2009) plus $\sim$15 further query groups
(phrase, variant, competition, Chinese-language layers) & zero hits for
$3\times n$ two-square tours; all future-work statements point elsewhere \\
OEIS \cite{oeis} & 50+ independent query strings over four rounds: sequence
variants (including $1,6,0,0,351,602,\dots$), keyword sweeps over all
knight-tour entries, comment sections of A169696 / A169764 / A070030 &
all zero; no deficient-board tour entry exists \\
formalisation ecosystem & Isabelle AFP \texttt{Knights\_Tour} (2022; boards
as arbitrary cell sets, no deficient-board theorem proved); Lean
\texttt{KnightMove} (2026, unpublished; vertex-deletion invariants, not
tour existence, no $3\times n$); Coq/Mizar; mathlib predicate layer &
no same-form result anywhere; mathlib has the Hamiltonian predicates but no
board object \\
recreational layer & Jelliss \emph{Knight's Tour Notes} index, $3\times n$
open/closed pages, holey-board pages & full-board theorems only; no
deficient $3\times n$ content \\
\hline
\end{tabular}%
}
\end{center}

The channels are named so a later reader can re-run them; if any channel
later returns a same-form record, the tier above weakens accordingly, and
this deposit will be amended by a later version rather than by editing this
record. Residual risks are ledgered in the C9 record (the Jelliss 2019
monograph and Chessics issues were checked at index level only; the Van
Rees classroom notes were reachable at bibliographic level only; the
StackExchange API was rate-limited during one round; Thai-language
literature was checked through the journal's own site).

\section{Scope and limitations}\label{sec:limits}

(a) The kernel results are Theorems~\ref{thm:open} and~\ref{thm:closed} and
nothing else. (b) The counting companion of the same family (the number of
open/closed tours of $K_n$) is auxiliary computational data of the source
task; it is not claimed here and not submitted to OEIS (the historical
OEIS sequence queries of the novelty record, all zero-hit, are included in
the record verbatim), and it is not part of this note. (c) The case $n=1$ of Theorem~\ref{thm:open} comes from the
pre-freeze correction F0b (recorded in the task card). (d) Scaffolding
disclosure, from the final adjudication: the cut-barrier section was
drafted by a pipeline worker run and integrated by the main agent (four
integration-level repairs); the final assembly section was written by the
main agent after a scratch pre-run. This affects provenance, not the
grading: the graded object is the final kernel-checked artifact. (e) The
delivered file relaxes elaboration resources
(\texttt{maxHeartbeats} 4{,}000{,}000 and \texttt{maxRecDepth} 8{,}000)
before the cut-barrier section; this affects elaboration time limits only,
not the kernel trust base. (f) An earlier search-based route to the
$n=6,7$ impossibility (a decidability checker) was abandoned after a
resource-exhaustion probe and is kept in the source task's archive; the
delivered route is the cut certificate. (g) The novelty search is scoped
to reachable channels; the residual risks listed in
Section~\ref{sec:novelty} are carried rather than resolved. (h) All listed
listings are byte-exact contiguous excerpts of the frozen artifacts
(LF-normalised where the source has CRLF-dominant mixed endings; the
byte-exact originals are shipped in \texttt{proofs/}).

\section*{Data availability}

The Lean sources (\texttt{proofs/}: frozen statement file, final proof
file, \texttt{lean-toolchain}), the verification and audit logs, and the
full novelty-search record are packed in the Zenodo deposit of this note
(DOI reserved at upload). The note is under CC BY 4.0, the code under MIT.
The upload steps are listed in \texttt{UPLOAD.md} in the source repository
(not shipped in this deposit).

\section*{Use of generative AI}

(a) AI performed: candidate generation and the four-round exclusivity
review; autoformalisation of the definitions and statements; the roundtrip
semantic check; proof search and proof-script writing (worker agents, with
a main-agent integration round on the cut-barrier section as recorded in
the ledger); error classification and repair; the computational exploration
scripts; and the assembly of this note including its verbatim listings.
(b) Model, node and budget: the proving phase ran on the node designated
by the author for each session (\texttt{anthropic/stepfun/step-5-preview}
after a mid-task switch directed by the author; \texttt{anthropic/a6api-main/kimi-k3}
before it), consuming $29$ agent-run samples of a permitted $36$; $4$
llm-calls of a permitted $6$ (one node smoke plus three roundtrip attempts
that failed at the transport layer, after which the semantic roundtrip was
re-routed through the agent channel and ledgered under the sampling
column); $49$ strict kernel compiles of a permitted $56$, plus
compile-debug probes ledgered separately; this note's lane consumed $0$
pipeline LLM calls against a cap of $4$ and three formal LaTeX rounds of a
permitted ten (two assembly rounds plus one audit-fix round), with three
compile-debug render probes and one kernel re-verification of the shipped
axioms listing ledgered separately; the assembly node was
\texttt{anthropic/a6api-main/deepseek-v4.1-flash}. (c) Final adjudication: all formal statements and proofs
reported in this note were checked by the Lean 4 kernel (Lean v4.34.0,
mathlib v4.34.0): every proof compiles with no \texttt{sorry}, and
\texttt{\#print axioms} reports only \{propext, Quot.sound,
Classical.choice\}. (d) The AI system is not an author; the human author
reviewed the evidence and is responsible for all content.

\section*{Author contributions}

CRediT-style: the human author adjudicated topic selection, statement
semantics, sign-off and deposit decisions; the AI4Math pipeline (an LLM)
generated the candidate proposition, the proof searches, the computational
exploration scripts and the text drafts. The AI system is not listed as an
author.

\begin{thebibliography}{9}
% Each entry verified on 2026-09-30 against Crossref title/volume/page
% metadata, the zbMATH Open API record (an:1538.05160), and live OEIS
% fetches (A169696). Raw responses:
% audit/bib-verify/ in the source repository. No entry is quoted from memory.
\bibitem{lean4}
L.~de Moura and S.~Ullrich.
\newblock The Lean 4 theorem prover and programming language.
\newblock \emph{CADE-28}, 2021. \url{https://doi.org/10.1007/978-3-030-79876-5_37}

\bibitem{mathlib}
The mathlib Community.
\newblock The Lean mathematical library.
\newblock \emph{CPP 2020}, 2020. \url{https://doi.org/10.1145/3372885.3373824}

\bibitem{millerfarnsworth}
J.~Miller and D.~Farnsworth.
\newblock Knight's tours on $3\times n$ chessboards with a single square
removed.
\newblock \emph{Open J.\ Discrete Math.} 3(1):56--59, 2013.
\url{https://doi.org/10.4236/ojdm.2013.31012}

\bibitem{srichote}
W.~Srichote, R.~Boonklurb, T.~Kaewwannarat, and S.~Singhun.
\newblock Closed knight's tours on $4\times n$ chessboards with two squares
removed.
\newblock \emph{Thai J.\ Math.}, Special Issue: Annual Meeting in
Mathematics 2021, 64--81, 2022. (Zbl 1538.05160)

\bibitem{bi}
Z.~Bi, S.~Butler, A.~DeGraaf, and E.~Doebel.
\newblock Knight's tours on boards with odd dimensions.
\newblock \emph{Involve} 8(4):615--627, 2015.
\url{https://doi.org/10.2140/involve.2015.8.615}

\bibitem{demaio}
J.~DeMaio and T.~Hippchen.
\newblock Closed knight's tours with minimal square removal for all
rectangular boards.
\newblock \emph{Math.\ Mag.} 82(3):219--225, 2009.
\url{https://doi.org/10.1080/0025570X.2009.11953624}

\bibitem{schwenk}
A.~J.~Schwenk.
\newblock Which rectangular chessboards have a knight's tour?
\newblock \emph{Math.\ Mag.} 64(5):325--332, 1991.
\url{https://doi.org/10.1080/0025570X.1991.11977627}

\bibitem{oeis}
OEIS Foundation Inc.
\newblock The On-Line Encyclopedia of Integer Sequences.
\newblock \url{https://oeis.org}

\bibitem{a169696}
OEIS Foundation Inc.
\newblock Sequence A169696: Number of undirected knight's tours on a $3\times n$
board.
\newblock \url{https://oeis.org/A169696}
\end{thebibliography}

\end{document}
